% page 44 du fac-simile (image p-043 du scan)
% bloc 165.1 x 242.0 mm ; marge gauche 8.0 mm
\pgc{-12.700mm}{0.000mm}{
\l{\cel{5}36}
\l{}
\l{\sou{apta}. Qua havas la karakteri necesa por (agar, facar, efikar).}
\l{\cel{3}- EFIS.}
\l{}
\l{\sou{apud}.\cel{2}Tre proxime, ma ne tam grande kam indikas \textquotedbl{}an\textquotedbl{}. - IL.}
\l{}
\l{\sou{apuntar}. (\sou{fusilo, kanono, lorno, teleskopo}). Direktar, mantenar}
\l{\cel{3}fixigita vers to quon onu vizas. - DEFIS.}
\l{}
\l{\sou{aquo}. (\sou{kemio}) Liquido qua konsistas ye du parti de hidrogeno}
\l{\cel{3}ed un parto de oxigeno.\cel{2}H2 O. -\cel{4}DEFIRS.}
\l{}
\l{\sou{aquaforto}.\cel{2}Azot-acido (AzO3H)dilutita per aquo, quan la grab-}
\l{\cel{3}isto uzas por korodar la kupro-plako. - DEFIRS.}
\l{}
\l{\sou{aquamarino}. Varietato di smeraldo, verda quale la mar-aquo.}
\l{\cel{3}- DEFIRS.}
\l{}
\l{\sou{aquarelo}.\cel{2}Pikturo lejera, en preske omna kazi sur papero,}
\l{\cel{3}facita ek farbi diafana, dilutita per aquo.- DEFIRS.}
\l{}
\l{\sou{aquario}. Tanko en qua onu entratenas animali o planti aquala.}
\l{\cel{3}- DEFIRS.}
\l{}
\l{\sou{aquatinto}. Graburo per aquaforto, qua imitas aquopikturo.-DEFIRS.}
\l{}
\l{\sou{aquedukto}.\cel{2}Tubego subsula od elevita sur arki, qua adduktas}
\l{\cel{3}aquo de loko ad altra. - DEFIRS.}
\l{}
\l{\sou{aquilegio}.\cel{2}(\sou{bot}.) Planto renunkulacea. - L.\cel{1}\sou{aquilegia}. - L.}
\l{}
\l{\sou{aquilono}.\cel{2}I.Vento de nordo. - II. (\sou{Literaturo)}.\cel{2}Vento kolda}
\l{\cel{3}e violentoza. - DEFIRS.}
\l{}
\l{\sou{aquirar}.\cel{2}(\sou{trans}.) (\sou{per, po}). Prokurar a su la posedo di ulo.}
\l{\cel{3}- DEFIS.}
\l{}
\l{\sou{-aro,}\cel{2}Sufixo qua, soldita a radiko nomala, indikas \textquotedbl{}la maxim}
\l{\cel{3}extensita) kolektajo, ensemblo\textquotedbl{} ek la kozi o personi nomata}
\l{\cel{3}da la radiko.}
\l{}
\l{\sou{-ar}. Finalo di la verbi ye infinitivo di prezento e di neprec-}
\l{\cel{3}izesoe pokala.}
\l{}
\l{\sou{arabesko}. Ornamento di pikturo o di skulturo, facita ek planti,}
\l{\cel{3}animali, fantazio-figuri interplektita kapricoze. - DEFIRS.}
\l{}
\l{\sou{arachar}.\cel{2}(\sou{trans., de, ek}). Deprenar (per violento, koakto, de}
\l{\cel{3}ulu to quon lu retenas. (\sou{Anke metaf.})\cel{2}- FS.}
\l{}
\l{\sou{arako}. Liquoro alkoholoza quan onu obtenas de rizo o de la}
\l{\cel{3}suko di kokoso. - DEFIRS.}
\l{}
\l{\sou{araknoido}. (\sou{anat.)}\cel{2}Un ek la tri envelopili di encefalo, mem-}
\l{\cel{3}brano seroza tre dina, situita inter du altra. - DEF.}
}
